\draw[separator] (11.85,6.15) -- (18.3,6.15);
\panelheading{12.0}{6.52}{(d) Segment-projection feedback}

\node[smallnote] at (15.10,7.00)
    {Reference segment $n=3$, direction $P=Z_2$};
\node[smallnote,anchor=west,text=ink] at (12.15,7.53)
    {Hermitian-matrix space};

% A plane schematically represents the six-dimensional local subspace.
\filldraw[draw=localblue!45,fill=localblue!7,line width=0.5pt]
    (12.35,8.95) -- (14.65,8.95) -- (15.10,9.35)
    -- (12.80,9.35) -- cycle;
\coordinate (localOrigin) at (12.85,9.15);
\coordinate (localFoot) at (14.15,9.15);
\coordinate (transportedDirection) at (14.15,7.98);
\draw[flow,draw=localblue,line width=0.9pt]
    (localOrigin) -- (transportedDirection);
\node[text=localblue,anchor=south east,inner sep=2pt]
    at ($(localOrigin)!0.5!(transportedDirection)$)
    {$V_3 Z_2 V_3^\dagger$};
\draw[detail] (transportedDirection) -- (localFoot);
\draw[flow,draw=localblue,line width=1.05pt] (localOrigin) -- (localFoot);
\draw[draw=muted,line width=0.45pt]
    (14.02,9.15) -- (14.02,9.02) -- (14.15,9.02);
\fill[ink] (localOrigin) circle[radius=0.025cm];
\fill[localblue] (localFoot) circle[radius=0.03cm];
\node[text=localblue,anchor=west,inner sep=1pt]
    at (14.25,9.15) {$A$};
% The commutator map connects the two spaces, without another projection.
\node[text=localblue,inner sep=1pt] (commutatorGenerator)
    at (13.35,10.65) {$\mathcal H_A$};
\node[text=localblue,anchor=base east,inner sep=1pt] (commutatorInput)
    at ([xshift=-0.45cm]commutatorGenerator.base west) {$A$};
\node[text=localblue,anchor=base,inner sep=0pt]
    at ($(commutatorInput.base east)!0.5!(commutatorGenerator.base west)$)
    {$\mapsto$};
\node[text=localblue,anchor=base west,inner sep=0pt]
    at ([xshift=1pt]commutatorGenerator.base east) {$\,=-i[A,\,\cdot\,]$};
\draw[flow,draw=localblue] (localFoot) -- (14.15,10.20)
    -| (commutatorInput.north);
\node[smallnote,text=localblue,fill=white,inner sep=3pt]
    at (13.725,9.75) {Local subspace (6 dimensions)};

\begin{scope}[shift={(0,0.80)}]
\node[smallnote,anchor=west,text=ink] at (12.15,10.65)
    {Superoperator space};

% A short line depicts the one-dimensional real span of the superoperator.
\coordinate (projectionOrigin) at (12.65,12.70);
\coordinate (projectionResidual) at (13.12,11.12);
\coordinate (projectionDirection) at (14.65,11.65);
\coordinate (projectionFoot) at
    ($(projectionOrigin)!(projectionResidual)!(projectionDirection)$);
\draw[draw=localblue!65,line width=0.65pt]
    (12.48,12.7893) -- (projectionDirection);
\node[smallnote,text=localblue,inner sep=1pt] (spanGenerator)
    at (15.02,11.12) {$\mathcal H_A$};
\node[smallnote,text=localblue,anchor=base east,inner sep=0pt]
    at (spanGenerator.base west) {$\operatorname{span}\{$};
\node[smallnote,text=localblue,anchor=base west,inner sep=0pt]
    at (spanGenerator.base east) {$\}$};
\node[smallnote,text=localblue] at (14.75,11.48) {(1 dimension)};
\draw[flow,draw=ink,line width=0.9pt]
    (projectionOrigin) -- (projectionResidual);
\node[anchor=east,inner sep=2pt]
    at ($(projectionOrigin)!0.5!(projectionResidual)$)
    {$\mathcal G$};
\draw[detail] (projectionResidual) -- (projectionFoot);
\draw[feedbackflow,line width=1.15pt]
    (projectionOrigin) -- (projectionFoot);
\fill[ink] (projectionOrigin) circle[radius=0.025cm];
\fill[feedbackgreen] (projectionFoot) circle[radius=0.03cm];
\coordinate (projectionAlong) at
    ($(projectionFoot)!0.12cm!(projectionOrigin)$);
\coordinate (projectionNormal) at
    ($(projectionFoot)!0.12cm!(projectionResidual)$);
\draw[draw=muted,line width=0.45pt]
    (projectionAlong) --
    ($(projectionAlong)+(projectionNormal)-(projectionFoot)$) --
    (projectionNormal);
\node[text=feedbackgreen,inner sep=1pt] (projectedMatrix)
    at (13.55,12.95) {$\lambda\mathcal H_A$};
\node[text=feedbackgreen,inner sep=2pt] (segmentFeature)
    at (16.45,12.45) {$\eta_{3,Z_2}=7\lambda/\sigma_\delta$};
\draw[feedbackflow] (projectedMatrix.east) -- (14.75,12.95)
    -- (14.75,12.45) -- (segmentFeature.west);

% Seven rows identify segments; two columns identify qubits.
\node[smallnote,text=ink] at (17.33,7.70) {Segment features};
\node[smallnote] at (16.35,8.13) {$n$};
\node[smallnote] at (17.06,8.13) {$q_1$};
\node[smallnote] at (17.60,8.13) {$q_2$};
\foreach \r in {0,...,6} {
    \pgfmathsetmacro{\y}{8.40+0.48*\r}
    \node[smallnote] at (16.35,{\y+0.20})
        {$\the\numexpr\r+1\relax$};
    \foreach \c in {0,1} {
        \pgfmathsetmacro{\x}{16.85+0.54*\c}
        \draw[draw=feedbackgreen!65,fill=feedbackgreen!10,line width=0.4pt]
            (\x,\y) rectangle ++(0.42,0.40);
    }
}
\draw[draw=feedbackgreen!65,fill=feedbackgreen!65,line width=0.4pt]
    (17.39,9.36) rectangle ++(0.42,0.40);
\node[smallnote] at (17.33,11.98) {$14$ features};
\draw[feedbackflow] (segmentFeature.east) -- (18.10,12.45)
    -- (18.10,9.56) -- (17.85,9.56);
\end{scope}
\draw[flow,draw=localblue] (commutatorGenerator.south) -- (13.35,11.05)
    -| (spanGenerator.north);
